% !TEX root = ../main.tex
\section{实验}
\label{sec:experiments}

\subsection{实验设置}
\label{sec:exp:setup}

\planned{$\tau^2$-bench 的航空与零售领域~\citep{barres2025tau2}；Qwen3-32B 和 Qwen3.5-9B~\citep{qwen2025qwen3,qwen2026qwen35}；启用批次不变内核的 vLLM~\citep{kwon2023vllm}；基线为串行执行与 \spork，后者采用与 \method 相同的候选隔离执行方式；明确披露历史策略。}

\subsection{RQ1：端到端延迟}
\label{sec:exp:latency}

\planned{逐任务配对延迟、P50 和 P95，以及任务级自助法（bootstrap）区间。}

\subsection{RQ2：无损性与任务质量}
\label{sec:exp:lossless}

\planned{轨迹完全一致性、最终状态相等性、奖励一致性，以及提交审计。}

\subsection{RQ3：机制}
\label{sec:exp:mechanism}

\planned{启动、接纳与拒绝统计；消融设置：延迟执行后继工具、后继轮次不探测、后继轮次仅探测一次，以及串行对照。}
